%% Career Ops Reiterated — ATS-Optimised LaTeX CV Template
%% Engine: XeLaTeX (recommended) or pdfLaTeX
%%
%% USAGE:
%%   1. Copy this file to cv/cv.tex
%%   2. Replace all placeholder values (search for YOUR_)
%%   3. Compile: python cv/generate_pdf.py --company "Acme" --role "AI Architect"
%%
%% ATS NOTES:
%%   - Single-column layout maximises ATS parse accuracy
%%   - No tables for the main body — plain text sections
%%   - Standard section headings: Experience, Education, Skills
%%   - UTF-8 throughout (xelatex handles special chars natively)
%%
%% TAILORED BUILDS:
%%   generate_pdf.py injects \renewcommand{\TargetCompany}{...} after
%%   \begin{document}, overriding the \providecommand defaults below.

\documentclass[11pt, a4paper]{article}

% ── Fonts & encoding ──────────────────────────────────────────────────────────
\usepackage[T1]{fontenc}
\usepackage[utf8]{inputenc}
% For XeLaTeX with system fonts, uncomment:
% \usepackage{fontspec}
% \setmainfont{Inter}[BoldFont=Inter Bold, ItalicFont=Inter Italic]

% ── Layout ────────────────────────────────────────────────────────────────────
\usepackage[top=18mm, bottom=18mm, left=20mm, right=20mm]{geometry}
\usepackage{parskip}
\usepackage{microtype}
\usepackage{enumitem}
\usepackage{titlesec}
\usepackage{xcolor}
\usepackage{hyperref}

% ── Colors ────────────────────────────────────────────────────────────────────
\definecolor{accent}{HTML}{1a73e8}   % change to your preferred accent color
\definecolor{muted}{HTML}{6b7280}

% ── Hyperlinks ────────────────────────────────────────────────────────────────
\hypersetup{colorlinks=true, urlcolor=accent, linkcolor=accent, hidelinks}

% ── Section headings ──────────────────────────────────────────────────────────
\titleformat{\section}
  {\large\bfseries\color{accent}}{}{0em}{}[\color{muted}\titlerule]
\titlespacing*{\section}{0pt}{10pt}{4pt}

% ── Bullet style ──────────────────────────────────────────────────────────────
\setlist[itemize]{leftmargin=1.4em, topsep=2pt, itemsep=1pt, parsep=0pt}

% ── Job-specific variables (overridden by generate_pdf.py) ────────────────────
\providecommand{\TargetCompany}{General Application}
\providecommand{\TargetRole}{YOUR_TARGET_ROLE}
\providecommand{\JobKeywords}{}

% ── Helpers ───────────────────────────────────────────────────────────────────
\newcommand{\cventry}[4]{%    {Company}{Role}{Dates}{Location}
  \textbf{#1} \hfill \textit{\small #3} \\
  \textit{#2} \hfill \textit{\small\color{muted} #4}
  \vspace{1pt}
}
\newcommand{\cveducation}[4]{%    {Institution}{Degree}{Dates}{Location}
  \textbf{#1} \hfill \textit{\small #3} \\
  \textit{#2} \hfill \textit{\small\color{muted} #4}
  \vspace{1pt}
}

% ═════════════════════════════════════════════════════════════════════════════
\begin{document}
\pagestyle{empty}

% ── Header ────────────────────────────────────────────────────────────────────
\begin{center}
  {\LARGE\bfseries YOUR\_FULL\_NAME}\\[4pt]
  \small
  YOUR\_CITY, YOUR\_COUNTRY
  $\cdot$ \href{mailto:your@email.com}{your@email.com}
  $\cdot$ +XX XXX XXXXXXX \\
  \href{https://linkedin.com/in/yourprofile}{linkedin.com/in/yourprofile}
  $\cdot$ \href{https://github.com/yourusername}{github.com/yourusername}
\end{center}

\vspace{4pt}

% ── Summary ───────────────────────────────────────────────────────────────────
\section{Summary}

% 2–3 sentences. Lead with your title + years of experience.
% End with what you're looking for or what makes you distinctive.
YOUR\_PROFESSIONAL\_SUMMARY

% ── Experience ────────────────────────────────────────────────────────────────
\section{Experience}

\cventry
  {Company Name}{Job Title}{Mon YYYY -- Present}{City, Country}
\begin{itemize}
  \item \textbf{[Quantified achievement]} — action verb + metric + context.
  \item Led cross-functional team of \textbf{N engineers} to deliver [outcome].
  \item Reduced [metric] from \textbf{X} to \textbf{Y} by [specific action].
  \item [Next achievement — action verb + metric + context]
\end{itemize}

\vspace{6pt}

\cventry
  {Previous Company}{Job Title}{Mon YYYY -- Mon YYYY}{City, Country}
\begin{itemize}
  \item [Achievement with metric]
  \item [Achievement with metric]
  \item [Achievement with metric]
\end{itemize}

\vspace{6pt}

\cventry
  {Earlier Company}{Job Title}{Mon YYYY -- Mon YYYY}{City, Country}
\begin{itemize}
  \item [Achievement with metric]
  \item [Achievement with metric]
\end{itemize}

% ── Projects / Portfolio ──────────────────────────────────────────────────────
\section{Projects}

\textbf{Project Name} \hfill
\href{https://github.com/yourusername/project}{github.com/yourusername/project}\\
One-line description. Hero metric: \textbf{KEY\_RESULT}. Stack: Python, FastAPI, PostgreSQL.

\vspace{4pt}

\textbf{Another Project} \hfill
\href{https://github.com/yourusername/project2}{github.com/yourusername/project2}\\
One-line description. Hero metric: \textbf{KEY\_RESULT}. Stack: Go, Kubernetes, Kafka.

% ── Skills ────────────────────────────────────────────────────────────────────
\section{Skills}

\begin{itemize}[leftmargin=0pt, label={}]
  \item \textbf{AI / ML:}          LLM APIs, RAG, vector databases, evals, MLflow
  \item \textbf{Cloud:}            AWS (EKS, Lambda, RDS), Azure, GCP, Terraform, Kubernetes
  \item \textbf{Languages:}        Python, Go, TypeScript
  \item \textbf{Architecture:}     Distributed systems, event-driven (Kafka), microservices, gRPC
  \item \textbf{Leadership:}       Cross-functional teams, technical roadmaps, OKR planning
\end{itemize}

% ── Education ─────────────────────────────────────────────────────────────────
\section{Education}

\cveducation
  {University Name}{B.Sc. / M.Sc. in Computer Science}{YYYY}{City, Country}

% ── Languages ─────────────────────────────────────────────────────────────────
\section{Languages}

English (fluent) \quad German (professional working proficiency) \quad [Your native language] (native)

\end{document}
